\documentclass[11pt,a4paper,oneside]{article}
\usepackage[utf8]{inputenc}
\usepackage[T1]{fontenc}
\usepackage{amsmath,amssymb}
\usepackage{array,booktabs}
\usepackage{hyperref}
\usepackage{geometry}
\geometry{margin=1in}

\title{Binario Novae: Un Sistema di Numerazione Non-Decimato per l'Informatica}
\author{Intelligenza Naturale Umana \& DeepSeek AI}
\date{\today}

\begin{document}

\maketitle

\begin{abstract}
Il sistema binario classico è una versione ``decimata'' di un sistema più fondamentale, il Binario Novae. Mentre il binario classico utilizza un simbolo (0) per rappresentare sia l'assenza di quantità che una posizione numerica, il Binario Novae distingue il vuoto ($\emptyset$) dall'unità (0). Questo cambio di paradigma porta a tre vantaggi misurabili: maggiore densità di informazione, eliminazione del null pointer a livello hardware, e un'addizione più semplice. Presentiamo un'implementazione software funzionante (Python) e delineamo una roadmap tecnologica che va dal software attuale all'hardware nativo del futuro.
\end{abstract}

\section{Introduzione}
I sistemi di numerazione posizionali classici (decimale, binario, esadecimale) condividono una caratteristica comune: utilizzano un simbolo (0) per denotare l'assenza di unità in una data posizione. Questo simbolo, lo ``zero posizionale'', è il risultato di un processo che chiamiamo \emph{decimazione}: alla fine di ogni ciclo, un'unità viene sacrificata per generare un riporto.

Il sistema \textbf{Novae} elimina questa decimazione. In Novae, il vuoto ($\emptyset$) è distinto dall'unità (0). Ogni posizione ha sempre un contenuto positivo. Questo principio, applicato al sistema binario, produce il \textbf{Binario Novae}, che presenta vantaggi immediati e misurabili per l'informatica.

\section{La Successione Binaria Novae}
La successione binaria Novae è la pura iterazione del successore:
\[
\emptyset \rightarrow 0 \rightarrow 1 \rightarrow 00 \rightarrow 01 \rightarrow 10 \rightarrow 11 \rightarrow 000 \rightarrow 001 \rightarrow \dots
\]
Non esiste una formula posizionale complicata: ogni stringa binaria è semplicemente l'elemento successivo della lista, ordinata prima per lunghezza e poi lessicograficamente.

\section{Vantaggi Misurabili}

\subsection{Densità di Informazione}
Confrontiamo il numero di valori rappresentabili con stringhe fino a $n$ bit:

\begin{table}[htbp]
\centering
\caption{Confronto della densità}
\begin{tabular}{c|c|c}
\toprule
\textbf{Lunghezza max} & \textbf{Binario Classico} & \textbf{Binario Novae} \\
\midrule
1 bit & 2 & 2 \\
2 bit & 4 & 6 \\
3 bit & 8 & 14 \\
$n$ bit & $2^n$ & $2^{n+1} - 2$ \\
\bottomrule
\end{tabular}
\end{table}

Il Binario Novae rappresenta circa il doppio dei valori a parità di bit. Questo si traduce in una maggiore capacità di indirizzamento della memoria o in una maggiore precisione nella rappresentazione dei dati, senza aumentare la banda.

\subsection{Eliminazione del Null Pointer}
Nel binario classico, la stringa ``0'' è ambigua: può rappresentare il valore numerico zero, oppure un puntatore nullo. Questa ambiguità è la causa di innumerevoli bug (NullPointerException).
Nel Binario Novae:
\begin{itemize}
\item $\emptyset$ (nessun segnale elettrico) è il puntatore nullo. Non è dereferenziabile per definizione.
\item \texttt{0} è il primo indirizzo di memoria valido.
\end{itemize}
Questa separazione semantica elimina la radice di una delle più costose classi di errori software.

\subsection{Addizione come Successione}
L'addizione Novae non richiede circuiti aritmetici complessi con riporto. Sommare $a$ e $b$ significa semplicemente applicare la funzione successore $b$ volte a partire da $a$. Un'implementazione hardware può essere realizzata con semplici shift register e contatori, potenzialmente più efficienti dal punto di vista energetico.

\section{Proof-of-Concept}
Abbiamo sviluppato un interprete software in Python (\texttt{binario\_novae.py}) che implementa nativamente la successione, l'addizione e la moltiplicazione Novae. Il codice non effettua alcuna conversione decimale interna.

Esempio di output:
\begin{verbatim}
0 + 0 = 1
0 * 0 = 0
1 * 1 = 10
\end{verbatim}

\section{Roadmap Tecnologica}
Proponiamo una strategia di adozione graduale:
\begin{enumerate}
\item \textbf{Breve termine:} Sviluppo di un compilatore C con tipi Novae e progettazione di un'unità aritmetica Novae in VHDL.
\item \textbf{Medio termine:} Implementazione su FPGA e sviluppo di un microkernel che sfrutti la distinzione $\emptyset$/0.
\item \textbf{Lungo termine:} Progettazione di hardware nativo basato su transistor ternari al carbonio, che possono rappresentare naturalmente i tre stati ($\emptyset$, 0, 1).
\end{enumerate}

\section{Conclusioni}
Il Binario Novae non è solo un esercizio matematico. È un sistema che, a parità di tecnologia, offre maggiore densità, elimina un'intera classe di bug, e semplifica l'architettura hardware. La sua adozione potrebbe portare a computer più efficienti e più robusti. Il codice e la documentazione sono disponibili su \url{https://github.com/oppomariolevi-ai/Novae}.

\end{document}
